% Compile with: uplatex file.tex && dvipdfmx file.dvi
%           or: latexmk -uplatex file.tex
%\documentclass[uplatex,dvipdfmx]{jsarticle}
\documentclass[dvipdfmx,a4paper,11pt]{jsarticle}
\usepackage{amsmath, amssymb}
\usepackage{tikz}

\title{scenario\_spin\_chain\_spectrum}
\date{\today}

\begin{document}
\maketitle

イジング模型（2サイト）でスピン演算子 $S^x_1$ の交換子 $[H,S^x_1]$ を求める。$S^x$ がどのように時間発展するか（運動方程式）の第1段。

イジング模型 2サイトで $S^x_1$ の交換子を求めて

\begin{equation}
  [H, S^{x}_{1}] = i J S^{y}_{1} S^{z}_{2}
  \label{eq:E1}
\end{equation}

同じ演算子の二重交換子 $[H,[H,S^x_1]]$ を求める（1段目の結果がそのまま2段目の入力になり、加速度に相当する）。

イジング模型 2サイトで $S^x_1$ の二重交換子を求めて

\begin{equation}
  [H, [H, S^{x}_{1}]] = J^{2} S^{x}_{1} \left(S^{z}_{2}\right)^{2}
  \label{eq:E2}
\end{equation}

無限温度でのモーメント $M_0=\langle(S^x)^2\rangle,\,M_2$ を求める。演算子の時間発展を Krylov 部分空間に展開する準備。

イジング模型 2サイトで $S^x_1$ の2次までのモーメントを求めて

\begin{equation}
  M_{0} = \frac{1}{4},\ M_{2} = \frac{1}{16} J^{2}
  \label{eq:E3}
\end{equation}

モーメントから、動的相関関数の Green 関数を表す連分数の係数 $b_1^2$ を求める（recursion 法）。

イジング模型 2サイトで $S^x_1$ の連分数を求めて

\begin{equation}
  b_{1}^2 = \frac{1}{4} J^{2}
  \label{eq:E4}
\end{equation}

連分数からスペクトル関数 $S(\omega)$ を求める。有限系なので離散的なデルタ関数の和になるはずである。

イジング模型 2サイトで $S^x_1$ のスペクトル関数を求めて

\begin{equation}
  G(z) = \cfrac{1}{z - \cfrac{\frac{1}{4} J^{2}}{z}},\quad S(\omega) = \frac{1}{2}\,\delta(\omega - \frac{1}{2} J) + \frac{1}{2}\,\delta(\omega + \frac{1}{2} J)
  \label{eq:E5}
\end{equation}

今度はハイゼンベルク模型・3サイト鎖の中心サイトでスペクトル関数を求める。係数が2つになり、中心ピークと側ピークが現れるはずである。

ハイゼンベルク模型 3サイト鎖で $S^x_2$ のスペクトル関数を求めて

\begin{equation}
  G(z) = \cfrac{1}{z - \cfrac{J^{2}}{z - \cfrac{\frac{5}{4} J^{2}}{z}}},\quad S(\omega) = \frac{5}{9}\,\delta(\omega) + \frac{2}{9}\,\delta(\omega - \frac{3}{2} J) + \frac{2}{9}\,\delta(\omega + \frac{3}{2} J)
  \label{eq:E6}
\end{equation}

格子の形を指定せず、配位数 $z$ による一般式としてハイゼンベルク模型の2次モーメントを求める。

ハイゼンベルク模型で一般の格子の $S^x$ の2次モーメントを求めて

\begin{equation}
  M_0 = \frac{1}{4},\ M_2 = \frac{1}{8} J^{2} z
  \label{eq:E7}
\end{equation}

ハイゼンベルク模型のエネルギーの高温展開を、格子の形を指定せず一般式（辺数・三角形数などのグラフ不変量）として4次まで求める。

ハイゼンベルク模型の高温展開を一般式で4次まで求めて

\begin{equation}
  \langle H^2 \rangle = \frac{3}{16} J^{2} |E|,\ \langle H^3 \rangle = \frac{9}{32} J^{3} N_\triangle - \frac{3}{32} J^{3} |E|,\ \langle H^4 \rangle = \frac{9}{32} J^{4} N_\square - \frac{9}{32} J^{4} N_\triangle - \frac{3}{32} J^{4} N_\wedge - \frac{3}{128} J^{4} |E| + \frac{27}{256} J^{4} |E|^{2}
  \label{eq:E8}
\end{equation}

一般式を具体的な正方格子 $3\times3$ に当てはめ、直前の一般式と整合するかを確認する。

ハイゼンベルク模型 正方格子 3x3 の高温展開を求めて

\begin{equation}
  \langle H^2 \rangle = \frac{9}{4} J^{2},\ \langle H^3 \rangle = -\frac{9}{8} J^{3} \quad (|E| = 12,\ N_\triangle = 0)
  \label{eq:E9}
\end{equation}

3頂点のグラフの三角形の数を直接数える（高温展開の3次の項に現れる三角形数 $N_\triangle$ を裏付ける確認）。

グラフ (1,2),(2,3),(1,3) の三角形の数を求めて

\begin{equation}
  \#\triangle = 1
  \label{eq:E10}
\end{equation}

有限クラスタの離散スペクトル（デルタ関数の和）と対比するため、無限系の連続スペクトル密度（平方根終端）を求める。

平方根終端の連続スペクトル密度を求めて

\begin{equation}
  t(z) = \frac{z - \sqrt{z^2 - 4 b_\infty^2}}{2},\quad S(\omega) = \frac{\sqrt{4 b_\infty^2 - \omega^2}}{2\pi b_\infty^2}\quad (|\omega| \le 2 b_\infty)
  \label{eq:E11}
\end{equation}

\end{document}